\section{Proof of the fixed-block theorem}\label{sec:fixed-block}

Fix an integer $k\ge2$. Put
\[
P_k(n)=\prod_{j=0}^{k-1}(n+j),\qquad
q_k(n)=\frac{P_k(n)}{\operatorname{rad}(P_k(n))}.
\]
The function $P_k$ is positive and strictly increasing on the positive
integers. Write $\omega(D)$ for the number of distinct prime divisors of
$D$. All implied constants in this section may depend on $k$.

Lemma~2.1 is a Pell-equation count; compare Lebowitz--Lockard
\cite[Lemma~8]{LL2025} and Tao \cite[Lemma~2.10]{Tao2026}.

\begin{lemma}\label{lem:support-fibre}
For every squarefree positive integer $Q$ and every real $X\ge2$,
\[
\#\{1\le z\le X:\operatorname{rad}(z(z+1))\mid Q\}
\ll 2^{\omega(Q)}\log(2X),
\]
with an absolute implied constant. Consequently, for every $\delta>0$,
every value of $\operatorname{rad}(P_k(n))$ has at most
\begin{equation}\label{eq:endpoint-fibre}
C_{k,\delta}X^{k\delta}\log(2X)
\end{equation}
preimages among the positive integers $n\le X$.
\end{lemma}

\begin{proof}
Write $z(z+1)=dy^2$, where $d$ is squarefree. Then $d\mid Q$ and
\[
(2z+1)^2-d(2y)^2=1.
\]
The case $d=1$ has no solution with $z\ge1$, since
$(2z+1-2y)(2z+1+2y)=1$ is impossible. For a fixed $d\ge2$, put
$A=2z+1$, $B=2y$ and $\alpha=A+B\sqrt d$. Thus
\[
1<\alpha<2A\le4X+2,\qquad \alpha^{-1}=A-B\sqrt d.
\]
If $\alpha_2>\alpha_1$ arise from two solutions, then
\[
\frac{\alpha_2}{\alpha_1}=U_0+V_0\sqrt d>1,
\qquad U_0,V_0\in\mathbb Z,\qquad U_0^2-dV_0^2=1.
\]
Its conjugate is its positive reciprocal. Hence
\[
U_0=\frac12\left(\frac{\alpha_2}{\alpha_1}
                 +\frac{\alpha_1}{\alpha_2}\right)\ge2,
\qquad
\frac{\alpha_2}{\alpha_1}\ge2+\sqrt3.
\]
There are $O(\log(2X))$ such solutions for each $d\mid Q$. Summing over
the squarefree divisors of $Q$ proves the first assertion.

For every $\delta>0$,
\[
2^{\omega(Q)}\ll_\delta Q^\delta.
\]
Indeed, $2\le p^\delta$ for all sufficiently large primes $p$, and the
remaining primes contribute a constant. If
$Q=\operatorname{rad}(P_k(n))$ with $n\le X$, then
$Q\le(X+k-1)^k\ll_k X^k$. Every preimage $m$ of $Q$ satisfies
$\operatorname{rad}(m(m+1))\mid Q$, since $k\ge2$. This gives
\eqref{eq:endpoint-fibre}.
\end{proof}

Erd\H{o}s and Straus \cite[p.~63]{ErdosStraus1977} also use a Pell equation
for products of consecutive integers.

For a positive integer $D$, let
\[
\rho_k(D)=\#\{x\bmod D:P_k(x)\equiv0\pmod D\}.
\]

\begin{lemma}\label{lem:block-roots}
For every positive integer $D$,
\[
\rho_k(D)\le(k-1)!\,k^{\omega(D)}.
\]
\end{lemma}

\begin{proof}
If $p>k-1$, the $k$ factors of $P_k$ have distinct roots modulo $p$.
At most one factor is divisible by $p$, so $\rho_k(p^j)=k$ for every
$j\ge1$.

For $p\le k-1$, put $C_p=v_p((k-1)!)$. Given an integer $x$, choose
$i\in\{0,\ldots,k-1\}$ for which $v_p(x+i)$ is maximal, taking
$v_p(0)=\infty$. For $j\ne i$, the ultrametric inequality and maximality
give
\[
v_p(x+j)\le v_p(j-i).
\]
Thus
\[
\sum_{j\ne i}v_p(x+j)
\le v_p\bigl(i!(k-1-i)!\bigr)\le C_p.
\]
If $p^j\mid P_k(x)$, then
$p^{\max(j-C_p,0)}\mid x+i$. For each $i$ there are at most $p^{C_p}$
such residue classes modulo $p^j$. Hence
$\rho_k(p^j)\le kp^{C_p}$. The Chinese remainder theorem now gives
\[
\rho_k(D)\le k^{\omega(D)}\prod_{p\le k-1}p^{C_p}
=(k-1)!\,k^{\omega(D)}.
\qedhere
\]
\end{proof}

Sanna \cite[Lemma~3.3]{Sanna2022} proves the related bound
$O_r(x/y^{1/2})$ for the number of $n\le x$ such that one of
$n,\ldots,n+r-1$ has powerful part exceeding $y$.

\begin{lemma}\label{lem:quotient-tail}
For all real $X\ge2$, $K\ge\max(k,2)$ and $0<\eta<1$,
\begin{equation}\label{eq:quotient-tail}
\#\{1\le n\le X:q_k(n)>K\}
\ll_{k,\eta}XK^{-1+\eta}
       +K^2(1+\log K)^{k-1}+X^{1/2}.
\end{equation}
\end{lemma}

\begin{proof}
Let $1\le n\le X$ satisfy $q_k(n)>K$. Suppose first that a prime $p>K$
divides $q_k(n)$. Then
$p^2\mid P_k(n)$. Since $p>K\ge k$, exactly one factor is divisible by
$p$, and $p^2\mid n+i$ for some $0\le i<k$. The number of these $n$ is
at most
\[
k\sum_{K<p\le\sqrt{X+k-1}}\left(\frac X{p^2}+1\right)
\ll_k\frac XK+X^{1/2}.
\]
The last bound follows by summing over all integers in place of primes.

Otherwise every prime factor of $q_k(n)$ is at most $K$. Multiply its
prime factors, with multiplicity, until the product first exceeds $K$.
This gives an integer $g$ such that
\[
g\mid q_k(n),\qquad K<g\le K^2,
\qquad g\operatorname{rad}(g)\mid P_k(n).
\]
For the last divisibility, if $p^a\mid g$, then
$a\le v_p(P_k(n))-1$. Lemma~\ref{lem:block-roots} bounds the number of
these $n$ by
\begin{equation}\label{eq:tail-crt}
(k-1)!\sum_{K<g\le K^2}k^{\omega(g)}
\left(\frac{X}{g\operatorname{rad}(g)}+1\right).
\end{equation}
For $0<\eta<1$, the Euler product
\[
\sum_{g\ge1}\frac{k^{\omega(g)}g^{1-\eta}}
                   {g\operatorname{rad}(g)}
=\prod_p\left(1+\frac{k}{p(p^\eta-1)}\right)
\]
converges, since its local factor is $1+O_{k,\eta}(p^{-1-\eta})$.
The first term of \eqref{eq:tail-crt} is therefore
$O_{k,\eta}(XK^{-1+\eta})$.

Let $\tau_k(g)$ count the ordered factorizations of $g$ into $k$
positive integers. Then $k^{\omega(g)}\le\tau_k(g)$ and, for $Y\ge1$,
\[
\sum_{g\le Y}\tau_k(g)
\le Y\left(\sum_{r\le Y}\frac1r\right)^{k-1}
\le Y(1+\log Y)^{k-1}.
\]
The second term of \eqref{eq:tail-crt} is
$O_k(K^2(1+\log K)^{k-1})$. This proves \eqref{eq:quotient-tail}.
\end{proof}

\begin{lemma}\label{lem:small-quotient}
For every real $X\ge2$ and $K\ge1$,
\[
\#\left\{(m,n):
\begin{gathered}
1\le m<n\le X,\quad q_k(m),q_k(n)\le K,\\
\operatorname{rad}(P_k(m))=\operatorname{rad}(P_k(n))
\end{gathered}\right\}
\ll_k K^2\log(2X).
\]
\end{lemma}

\begin{proof}
Put $\mu=q_k(m)$ and $\nu=q_k(n)$. Strict monotonicity of $P_k$ gives
\begin{equation}\label{eq:quotient-curve}
1\le\mu<\nu\le K,\qquad \nu P_k(m)=\mu P_k(n).
\end{equation}
Set
\[
\xi=2m+k-1,\qquad \zeta=2n+k-1,
\qquad
R_k(U)=\prod_{j=0}^{k-1}(U+2j-k+1).
\]
Then $R_k(\zeta)=2^kP_k(n)$. With
\[
\mathcal H_k=\{h:1\le h\le k-1,\ h\equiv k-1\pmod2\},
\qquad A_k=\sum_{h\in\mathcal H_k}h^2,
\]
we have
\[
R_k(U)=U^kG_k(U),\qquad
G_k(U)=\prod_{h\in\mathcal H_k}\left(1-\frac{h^2}{U^2}\right).
\]
For $U\ge k+1$, the function $G_k$ is strictly increasing and
\begin{equation}\label{eq:centered-defect}
0<G_k(U)<1,\qquad 1-G_k(U)\le\frac{A_k}{U^2}.
\end{equation}
Indeed, $\mathcal H_k$ is nonempty and each factor is positive and
strictly increasing. The last inequality follows from
$1-\prod_j(1-x_j)\le\sum_jx_j$ for $0\le x_j\le1$.

Fix $\mu,\nu$ in \eqref{eq:quotient-curve}, and put
$\gamma=(\nu/\mu)^{1/k}$ and $\lambda=\zeta/\xi$.
Thus $\zeta>\xi\ge k+1$ and
\[
\gamma^k=\lambda^k\frac{G_k(\zeta)}{G_k(\xi)},\qquad 1<\lambda<\gamma.
\]
Using \eqref{eq:centered-defect}, we obtain
\begin{align*}
0<\gamma-\lambda
&=\gamma\left(1-\left(\frac{G_k(\xi)}{G_k(\zeta)}\right)^{1/k}\right)\\
&\le\gamma\bigl(1-G_k(\xi)^{1/k}\bigr)\\
&\le\gamma\bigl(1-G_k(\xi)\bigr)\\
&\le\frac{A_k\gamma}{\xi^2}.
\end{align*}
Write $\lambda=h/q$ in lowest terms, with $h,q>0$. Since $q\mid\xi$,
\[
0<\gamma-\frac hq\le\frac{A_k\gamma}{q^2},\qquad q\le2X+k-1.
\]
In a block $2^j\le q<2^{j+1}$, these fractions lie in an interval of
length $A_k\gamma/2^{2j}$. Distinct reduced fractions in the block are
more than $2^{-2j-2}$ apart. There are at most $1+4A_k\gamma$ fractions
in each block, and hence $O_k(\gamma\log(2X))$ possible values of $\lambda$.

For each such $\lambda$, every possible $\xi$ is a root of
\[
\nu R_k(U)-\mu R_k(\lambda U)=0.
\]
This polynomial has degree $k$, since its leading coefficient is
$\nu-\mu\lambda^k=\mu(\gamma^k-\lambda^k)>0$.
It has at most $k$ real roots. Each root determines
at most one pair $(m,n)$. Therefore, for fixed $\mu,\nu$, the number of
pairs in \eqref{eq:quotient-curve} is
\[
\ll_k (\nu/\mu)^{1/k}\log(2X).
\]
Finally,
\[
\sum_{1\le\mu<\nu\le K}(\nu/\mu)^{1/k}
\le\sum_{\nu\le K}\nu^{1/k}\frac{k}{k-1}\nu^{1-1/k}
\ll_k K^2.
\qedhere
\]
\end{proof}

\begin{proof}[Proof of Theorem~\ref{thm:main}]
Let $\varepsilon>0$. A counted pair either has both quotients at most
$K$, or has an endpoint in the exceptional set of
Lemma~\ref{lem:quotient-tail}. Multiplying the latter count by
\eqref{eq:endpoint-fibre} and applying Lemma~\ref{lem:small-quotient}
gives, for $K\ge\max(k,2)$,
\begin{align}\label{eq:fixed-block-assembly}
F_{k,k}(X)\ll_{k,\delta,\eta}{}&K^2\log(2X)\notag\\
&+X^{k\delta}\log(2X)
\left(XK^{-1+\eta}+K^2(1+\log K)^{k-1}+X^{1/2}\right).
\end{align}
The fibre bound counts all possible partners $1\le m\le X$, so this
also covers pairs whose smaller endpoint is exceptional.

Ignoring the factors $X^{k\delta}$, $K^\eta$ and the logarithms,
$K=X^\tau$ gives the exponent $\max(2\tau,1-\tau,1/2)$, minimized at
$\tau=1/3$. Set $K=\max(k,2,X^{1/3})$. If $X\ge k^3$, then
$K=X^{1/3}$. With this
choice fixed, take
\[
\eta=\min\left(\frac12,\frac\varepsilon2\right),
\qquad \delta=\frac\varepsilon{4k}.
\]
Every power of $X$ in \eqref{eq:fixed-block-assembly} is at most
\[
\frac23+k\delta+\frac\eta3
\le\frac23+\frac{5\varepsilon}{12}.
\]
The logarithmic factors together are
$O_{k,\varepsilon}(X^{\varepsilon/2})$. Hence
\[
F_{k,k}(X)\ll_{k,\varepsilon}X^{2/3+\varepsilon}.
\]
The bounded range $2\le X<k^3$ is absorbed into the implied constant.
\end{proof}
